% Chapter 3: Related work. Draft for the lead, 23 Nov 2026. Every number is traced in ch03_related_NOTES.md.
% Citation keys: refs_extra.bib, ch02_background.bib, ch03_related.bib. The related-work tables tab:rw_main, tab:rw_abs,
% tab:rw_out and tab:rw_counts live in related_work.tex (appendix); they are referenced here, not rebuilt.
\chapter{Related work and why published numbers do not compare}\label{ch:related}

This chapter reviews the public EL data (Section~\ref{sec:rw_data}), the pixel-level numbers published on them
(Section~\ref{sec:rw_numbers}), the methods that our benchmark suite draws on (Section~\ref{sec:rw_methods}), work on label
noise and benchmark practice (Section~\ref{sec:rw_noise}), and the semi-supervised framing this project started from
(Section~\ref{sec:rw_ssl}). Its argument is that published EL segmentation numbers cannot be compared with each other or with
ours: almost none states its split and metric together, and none gives an interval on its headline number. The protocol of
Chapter~\ref{ch:protocol} is built to close exactly these gaps.

\section{Public EL datasets}\label{sec:rw_data}

Table~\ref{tab:rw_datasets} lists the public EL data with pixel or region labels that the published numbers rest on, and two
image-level or box-level sets that matter for provenance. Three facts about them shape everything later in this chapter.

\begin{table}[htbp]
\centering
\scriptsize
\caption{Public EL datasets relevant to this work. ``Leakage'' lists what is known about overlap between train and test;
``Use here'' says how this thesis uses the set.}
\label{tab:rw_datasets}
\begin{tabular}{>{\raggedright}p{2.4cm} >{\raggedright}p{2.9cm} >{\raggedright}p{2.6cm} >{\raggedright}p{3.6cm} >{\raggedright\arraybackslash}p{2.4cm}}
\toprule
Dataset & Labels & Size & Splits and leakage & Use here \\
\midrule
BenchmarkEL \texttt{20211104} \citep{pratt2023benchmark} & pixel masks, 24 classes & 582 cells: 478 / 54 / 50 &
test set hand-picked for defects; 8 test cells are the same physical cell as a pool cell (Table~\ref{tab:T2}) &
pool P532 and the legacy test set \\
BenchmarkEL \texttt{20220723}, \texttt{20221008} & pixel masks, 28 and 29 classes; \texttt{20221008} corrects many crack
labels (dotted labels joined, lines widened) & 642 and 695 cells &
splits reshuffled: 186 of the 582 cells change split between the first two releases; the test sets of the later releases
(66 and 72 cells) each hold 51 of our pool cells & source of our revised crack and inactive labels and of the E1 cells \\
E-SCDD \citep{hanifi2025escdd} & the \texttt{20221008} masks repackaged (crack and inactive pixels unchanged), one class
added; 117 newly labelled real cells; 93 synthetic cells & 903 files: 693 + 117 + 93 & 44 of its 82 test cells come from
modules that also have train cells & download route of the revised labels; synthetic files excluded \\
UCF-EL \citep{fioresi2022ucfel} & polygons; ``crack'' is a region class (about crack plus inactive), contact, interconnect,
corrosion & public release 9,562 / 2,289; paper data 17,064 incl.\ 256 simulated & 25 to 44\% of the public test cells have
a physical twin in train & one pre-protocol pretraining test \\
PV-Vision (LBNL) \citep{chen2023pvvision} & pixel masks: crack, cross, busbar, dark & 1,837 cells sampled from 29,664 cells
of 445 module images & random cell split; cells of one module can fall in train and test & none \\
MultiSolSegment \citep{sanghi2026multisolsegment} & multi-label masks: crack, busbar, dark, non-cell & 585 mono-Si images &
augmented before a random split, so flipped copies can sit on both sides & none \\
PVEL-S \citep{duan2025caak} & binary masks & 1,200 poly-Si cells from PVEL-AD; 840 / 360 & evaluated on the validation set
in \citet{duan2025caak} & none \\
ELPV \citep{deitsch2019elpv} & image-level defect probability & 2,624 cells & contains the BMRK cells of BenchmarkEL & none
(source of BMRK) \\
PVEL-AD \citep{su2023pvelad} & boxes & 36,543 NIR cells & by request & none \\
\bottomrule
\end{tabular}
\end{table}

First, the main pixel-labelled set exists in three releases with reshuffled splits. A number on ``BenchmarkEL'' therefore
names one of three test sets, and the later test sets each hold 51 cells of our training pool; a model trained on either later
release split would have trained on 38 of the 50 legacy test cells. Second, the defect vocabulary is not shared: UCF-EL's
``crack'' merges thin polygons around crack lines with outlines of darkened and isolated regions, so it corresponds roughly to
BenchmarkEL crack plus inactive; a perfect UCF-EL crack mask shifted by 1~px keeps IoU 0.954, against 0.559 on the
BenchmarkEL legacy test masks (Chapter~\ref{ch:background}). Third, where we could check, train and test sets share modules
or even physical cells. None of the pixel-labelled sets we could check comes with a module-disjoint split. This thesis therefore
tests on E1, 106 cells from modules not seen in training (unseen modules: supported, not proven; Chapter~\ref{ch:data}).

\section{Published numbers and why they cannot be compared}\label{sec:rw_numbers}

We collected the papers we could find that report a pixel-level EL segmentation number (or a pixel-level anomaly map) on
solar cells from 2018 to 2026, and recorded for each its data and release, split, classes, metric and averaging, numbers,
interval and code, with the page each value comes from. Twenty-two papers were read in full text; seven more are known only
from their abstracts or second-hand, and a few could not be accessed at all. The full tables are in Appendix~\ref{app:related}: Table~\ref{tab:rw_main} for the full-text papers
\citep{pratt2021unet,pratt2023benchmark,eesaar2023seipv,mazen2024segnet,torpey2024pretraining,jha2025pvs3,fioresi2022ucfel,%
hanifi2025multiclass,hanifi2025escdd,ramadan2026hybrid,lodhi2026efsegnet,chen2022pvsc,chen2023pvvision,sohail2023fault,%
elbahri2026autoencoder,sanghi2026multisolsegment,duan2025caak,gunasekara2026stcnet,saleh2026anomaly,lin2026eiae,su2021sigan,%
hijjawi2026instance}, Table~\ref{tab:rw_abs} for the abstract-level rows
\citep{mahboob2025segformer,nigin2025mcxnet,bramantyo2026fracture,liu2027pvread,zhou2024siif,wang2024rern,narayanan2024solar},
and Table~\ref{tab:rw_out} for work left out because it is image-level, box-only or not EL. Discovery used our paper ledger and
literature notes, a forward-citation walk, local libraries and web search; conference-only and non-English work is
under-covered.

\paragraph{What the record shows.} Table~\ref{tab:rw_counts} counts the problems over the 22 full-text papers (a paper counts
once, and a problem counts if any of its rows has it).
\begin{itemize}
  \item \textbf{Split.} 0 of 22 papers state a module- or panel-disjoint split. 4 do not state the split at all and 1 states
  it in two conflicting ways. Where the paper or our own check allowed it, test data overlap the training data or the test
  set selects the model in 7 of 22; for 14 others the overlap cannot be checked from the paper.
  \item \textbf{Uncertainty.} 18 of 22 report no interval, seed count or repeat on any number. The other 4 give an interval
  only on a side quantity or a subset; none gives one on its headline test number.
  \item \textbf{Metric.} 15 of 22 leave the metric undefined (averaging, absent-class rule, class set or background not
  stated) or print numbers that are internally inconsistent. 10 of 22 put background, or a no-defect class, inside the
  headline mean.
  \item \textbf{Release.} 7 of 22 do not state the release, or report on data that is not the public release.
  \item \textbf{What nobody measures.} 0 of 22 use a distance-tolerant metric for the 1 to 2~px crack labels, and 0 of 22
  report a number on unseen modules. Only 2 of 22 state release, split and metric together, and one of these two evaluates
  on its validation set \citep{duan2025caak}.
\end{itemize}

\paragraph{One data family, many test sets.} Only three papers score on the legacy test set of BenchmarkEL
\citep{pratt2023benchmark,eesaar2023seipv,jha2025pvs3}, the last with its split not stated. \citet{torpey2024pretraining} use
release \texttt{20220723}; \citet{mazen2024segnet}, \citet{ramadan2026hybrid} and \citet{hanifi2025multiclass} use
\texttt{20221008}, and the last re-split it without releasing the split. Yet \citet{mazen2024segnet} and
\citet{ramadan2026hybrid} rank their numbers against legacy test-set rows of other papers, a comparison across test sets, and
\citet{jha2025pvs3} set their pooled IoU against Pratt's per-image medians, a comparison across metrics. The same holds for
UCF-EL, which appears as the official release, as the papers' own 80/10/10 splits and as the unreleased paper data, and for
PVEL-S.

\paragraph{Test data that the model has seen.} The legacy test set was not drawn at random: its authors assigned ``fifty
images with a prevalence of cracks, gridline defects, and inactive areas'' to it \citep{pratt2023benchmark}. Our module audit
found that 8 of its cells are the same physical cell as a cell of the training pool (7 confirmed by pixels, 1 resting on the
cell-list link only), 7 more are neighbouring crops of pool images, and 19 more share a module with pool cells
(Table~\ref{tab:T2}). Elsewhere, the released split lists of \citet{jha2025pvs3} put 1,769 of the 2,289 public UCF-EL test
images in their training lists; \citet{fioresi2022ucfel} validate on the test set after each epoch; \citet{duan2025caak}
evaluate on their validation set; and \citet{sanghi2026multisolsegment} and \citet{saleh2026anomaly} describe augmenting the
images before splitting them.

\paragraph{Convention moves the number as much as the method.} The same predictions give very different numbers under the
conventions in use. We re-scored our own full-pool PF models (5 seeds, mean over seeds) on the 50 legacy test cells with the
release labels (in-distribution, known leakage, context only; Table~\ref{tab:T13b}). The mean of the three defect medians in
the per-image style of \citet{pratt2023benchmark} is 0.405 when cells without the class are left out (rule A), 0.202 when they
score 0 (rule B), and 0.541 when they score 1 if nothing is predicted (rule C). A per-image mean IoU in the style of
\citet{eesaar2023seipv} is 0.855 when absent and unpredicted classes score 1, and 0.545 when they are skipped; the pooled micro
IoU is 0.891. On what are presumably the same 50 cells (PV-S3 does not state its split), crack is reported as 0.25 by
\citet{pratt2023benchmark} (per-image median) and as 0.58 by \citet{jha2025pvs3} (pooled IoU); the two
differ in model, metric and possibly split, so the gap cannot be assigned to any one of them. The ``IoU'' 0.7447 of \citet{eesaar2023seipv} equals
Dice / (2 $-$ Dice) of its Dice 0.8531 within 0.001, so the two are one pooled pixel quantity dominated by background and
padding; its headline ``mIoU'' 0.8604 has no stated averaging, and \citet{ramadan2026hybrid} reprint it as a
frequency-weighted IoU. On PVEL-S, \citet{gunasekara2026stcnet} print a mean IoU of 95.98\% on the validation
set and 72.52\% on the 360 test images. The label release alone also moves levels: re-scoring
the same cross-validated models on each release, PF scores strict crack IoU 0.320 on the \texttt{20211104} labels and 0.365
on our label set v1.1, and relF1@2 0.822 and 0.846, while contrasts between methods move by at most 0.008
(Table~\ref{tab:T3c}).

\paragraph{The published rows we can relate to our data.} Table~\ref{tab:T13a} copies the four published rows that report on
the 50 legacy test cells: model 4c and the best per-class values of \citet{pratt2023benchmark}, SEiPV-Net
\citep{eesaar2023seipv}, and the fully supervised PV-S3 row on its first BenchmarkEL subset \citep{jha2025pvs3}. None is
comparable with our protocol as published. Pratt's absent-class rule is not stated and only inferred (rule A); SEiPV-Net's
averaging is not defined, so only a family of readings is possible; PV-S3 states neither its split nor its background rule;
all are single runs; and the 50 cells overlap our training pool. Four further papers
\citep{mazen2024segnet,ramadan2026hybrid,hanifi2025multiclass,torpey2024pretraining} report on other test sets that also
overlap our pool. Table~\ref{tab:T13b} re-scores our own models in each published convention, as context only; no row of
Tables~\ref{tab:T13a} to~\ref{tab:T13c} supports a claim to beat a published model.

In our judgement, a published EL number is comparable only if it states release, label set, split, class list, metric
convention, operating-point rule and seed count together. None of the 22 full-text papers states all of them; the protocol
of Chapter~\ref{ch:protocol} fixes all seven and requires every entry to state them.

\section{Methods in our benchmark suite}\label{sec:rw_methods}

The suite of Chapter~\ref{ch:methods} runs recent methods from outside the EL literature next to the models of the EL papers,
all under one protocol. Table~\ref{tab:rw_arms} gives the origin of each arm; the results are in Chapter~\ref{ch:methods} and
in the tables named in the last column. Every arm is an adaptation, not a copy: it runs under the PF recipe (20,000 steps,
batch 8, 512~px, label set v1.1) unless the table says otherwise, not under its authors' training recipe.

\begin{table}[htbp]
\centering
\scriptsize
\caption{Origins of the arms of the benchmark suite. PF is U-Net++ with an ImageNet ResNet-34 encoder and the
encoder-preserving fine-tune; ``Results'' names the tables that report each arm.}
\label{tab:rw_arms}
\begin{tabular}{>{\raggedright}p{3.4cm} >{\raggedright}p{4.1cm} >{\raggedright}p{4.0cm} >{\raggedright\arraybackslash}p{2.2cm}}
\toprule
Arm (name in tables) & Origin & What it changes & Results \\
\midrule
PF, HOUSE & U-Net++ decoder \citep{zhou2018unetpp,zhou2020unetpp}; ResNet-34 encoder \citep{he2015resnet};
\citep{iakubovskii2019smp} & baseline; PF freezes the encoder for 1,000 steps, then ramps it to 0.1 times
the learning rate over the next 1,000 steps, with no encoder weight decay & Tables~\ref{tab:T5}, \ref{tab:T6} \\
unet\_r34\_pf & U-Net \citep{ronneberger2015unet} & decoder & \ref{tab:T15}, \ref{tab:T16a} \\
dlv3p\_r34\_pf & DeepLabV3+ \citep{chen2018deeplabv3plus} & decoder & \ref{tab:T15}, \ref{tab:T16a} \\
HRNET\newline (hrnet\_w18\_unetpp\_pf) & HRNet \citep{wang2021hrnet} & encoder that keeps a high-resolution branch &
\ref{tab:T6}, \ref{tab:T16b} \\
segformer\_mitb2\_pf & SegFormer \citep{xie2021segformer} & transformer encoder (MiT-B2) and MLP decoder &
\ref{tab:T15}, \ref{tab:T16a} \\
hanifi\_vgg16\_unet\_pf & VGG-16 U-Net of \citet{hanifi2025multiclass} & published EL model & \ref{tab:T15},
\ref{tab:T16a} \\
ucf\_dlv3\_r50\_pf & DeepLabV3 / ResNet-50 of UCF-EL \citep{fioresi2022ucfel} & published EL model & \ref{tab:T15},
\ref{tab:T16a} \\
pratt\_dlv3p\_wcec & weighted cross-entropy of \citet{pratt2023benchmark} on DeepLabV3+ & published EL loss &
\ref{tab:T15}, \ref{tab:T16a} \\
skr01\_pf\_v11 & Skeleton Recall loss \citep{kirchhoff2024skeletonrecall} & loss term for thin structures &
\ref{tab:T15}, \ref{tab:T16a} \\
jml\_hard, soft\_band\_jml & Jaccard metric loss \citep{wang2023jml} & loss that accepts soft labels; soft band around
crack labels & \ref{tab:T15}, \ref{tab:T16a} \\
SWAD (swad\_const) & SWAD \citep{cha2021swad} & dense weight averaging over a fixed window of training & \ref{tab:T6},
\ref{tab:T16b} \\
ema\_wise\_posthoc & EMA of weights \citep{moralesbrotons2024ema} & exponential moving average of weights, read post hoc &
\ref{tab:T15}, \ref{tab:T16a} \\
encoder\_spd & Selective Projection Decay \citep{tian2024spd} & selective weight decay towards the pretrained encoder,
in place of PF's lower encoder learning rate & \ref{tab:T15}, \ref{tab:T16a} \\
rcs & rare class sampling of DAFormer \citep{hoyer2022daformer} & sampling of training cells by class frequency &
\ref{tab:T15}, \ref{tab:T16a} \\
gin\_ipa & GIN and IPA augmentation \citep{ouyang2021gin} & augmentation for single-source domain generalisation &
parked, not run (\ref{tab:T16a}) \\
D4 & none: 8-view flip and rotation test-time augmentation, cited to no publication & inference & \ref{tab:T7a},
\ref{tab:T16b} \\
msc\_tta & scale views motivated by \citet{huang2026pie} & extra test-time views on top of D4 & \ref{tab:T7b} \\
\bottomrule
\end{tabular}
\end{table}

Two groups of arms deserve a word. The first are the models of the EL papers re-run under equal training. In two-seed screens
on the earlier split cv5v2 (which carried a module leak, Chapter~\ref{ch:protocol}), the UCF-EL model scored relF1@2 $-$0.0773
[$-$0.0892, $-$0.0594] and Pratt's weighted cross-entropy $-$0.0932 [$-$0.1113, $-$0.0774] against PF, while the VGG-16 U-Net
of Hanifi et al.\ was not separated from PF ($-$0.0064 [$-$0.0189, +0.0042]) (Table~\ref{tab:T16a}). Under one protocol, the
published EL models sit at or below our U-Net++ baseline. The second are methods built for a different problem: SWAD and
GIN-IPA target domain generalisation \citep{cha2021swad,ouyang2021gin}, rare class sampling targets domain adaptation
\citep{hoyer2022daformer}, and Skeleton Recall targets thin tubular structures
\citep{kirchhoff2024skeletonrecall}. We tried them because the unseen-module problem of this thesis is a domain shift and the
crack is a thin structure, not because they were shown to work on EL.

\section{Label noise and benchmark practice}\label{sec:rw_noise}

Our audit of the labels (Chapter~\ref{ch:data}) draws on work about label errors and benchmark validity outside EL. Of these
works we read the abstracts and, for one, the method section. Confident learning estimates which labels are wrong from a
model's predicted probabilities \citep{northcutt2021confident}, and Northcutt and colleagues found label errors in the test sets
of ten common benchmarks, enough to change the ranking of models \citep{northcutt2021pervasive}. For segmentation, label quality
can be scored per image from the predictions of any trained model \citep{lad2023labelquality}, and label errors can be found
at the level of connected components \citep{rottmann2023labelerrors}. Near-duplicates between train and test sets, found
in CIFAR, bias comparisons towards memorisation \citep{barz2020cifair}; contradictory annotations of the same data set an
upper bound on any measured score \citep{tschirschwitz2025label}; and bias in visual datasets is surveyed by
\citet{fabbrizzi2022bias}. Each of these has an EL counterpart in our data: the label-error candidates and the double annotation of the pool, the physical-cell duplicates between the legacy test
set and the pool, and the image statistics that name the source (Chapter~\ref{ch:background}). Chapter~\ref{ch:data} measures
each of them.

\section{Semi-supervised learning: why the framing was dropped}\label{sec:rw_ssl}

This project began as semi-supervised segmentation: 11,938 unlabelled CSIR and ARTS cells were available on our servers,
from a public pool of about 159,000 cells of the same two sources. The literature gives a reason to expect little from this
with all labels in hand: the gain from unlabelled images shrinks as the labelled share grows. On the Pascal benchmark of
\citet{yang2025unimatchv2}, for example, the labelled-only baseline with a DINOv2-S encoder scores 67.0\% mIoU with 1/16 of the
labels and 83.7\% with half of them, while UniMatch V2 scores 79.0\% and 86.7\%. On EL images, \citet{torpey2024pretraining}
compare pretraining regimes on release \texttt{20220723}; their two semi-supervised rows were run once each, against five runs per
regime for the others. \citet{jha2025pvs3} report semi-supervised results on BenchmarkEL subsets and on UCF-EL, the latter with
test images inside their released training lists (Section~\ref{sec:rw_numbers}).

Our own pre-protocol results agree with this shape, on the legacy split and the 50 legacy test cells (in-distribution, known
leakage, context only; Table~\ref{tab:T21}). At 5\% labels, pseudo-labelling with separate batch normalisation and CutMix
improved inactive IoU by +0.117 [+0.042, +0.189] and crack relF1@2 by +0.033 [+0.006, +0.055]; from 20\% of the labels upward
no difference was detectable, and adding the external unlabelled pool at full labels gave no detectable gain. A null of this kind closes
only the recipes that were run, under their conditions; it does not show that unlabelled data is useless for EL. In our
reading it did show that, with every label in use, the measurable problems of this work lay elsewhere: in the split, the
metric and the unseen modules. The project then dropped the semi-supervised framing and became a benchmark, for the
reasons of Section~\ref{sec:rw_numbers} (Chapter~\ref{ch:intro}); Chapter~\ref{ch:legacy} gives the history.
